% Generated by `taskgen course examples` from modular.ntt@mersenne (seed 1).
% Do not edit: change examples.yaml and regenerate.
\begin{example}[ТЧП Мерсенна, $m = 31$, $n = 5$]\label{ex:ntt-mersenne}
Вычислить циклическую свёртку последовательностей \[a = (0, 5, 4, 4, 9), \qquad b = (9, 8, 4, 9, 3)\] длины $n = 5$ в кольце $\mathbb{Z}/31\mathbb{Z}$, $31 = 2^{5} - 1$, с помощью теоретико-числового преобразования с корнем $\omega = 2$ из единицы степени $5$. Проверить ответ прямым вычислением свёртки.

\par\noindent\textit{Решение.} $31 = 2^{5} - 1$, поэтому $2^{5} \equiv 1$, и $\omega = 2$ имеет порядок $5 = n$ (5 --- простое число); умножение на $2^e$ --- циклический сдвиг $5$-битового слова. Прямое преобразование: $A_k = \sum_j a_j \omega^{jk}$, матрица $(\omega^{jk} \bmod m)$:\[W = \begin{pmatrix}1 & 1 & 1 & 1 & 1\\1 & 2 & 4 & 8 & 16\\1 & 4 & 16 & 2 & 8\\1 & 8 & 2 & 16 & 4\\1 & 16 & 8 & 4 & 2\end{pmatrix}.\]\[\begin{array}{c|l|r|r}k & A_k = \sum_j \omega^{jk} a_j & & A_k \bmod 31 \\ \hline 0 & 0 + 5 + 4 + 4 + 9 & 22 & 22 \\ 1 & 0 + 2 \cdot 5 + 4 \cdot 4 + 8 \cdot 4 + 16 \cdot 9 & 202 & 16 \\ 2 & 0 + 4 \cdot 5 + 16 \cdot 4 + 2 \cdot 4 + 8 \cdot 9 & 164 & 9 \\ 3 & 0 + 8 \cdot 5 + 2 \cdot 4 + 16 \cdot 4 + 4 \cdot 9 & 148 & 24 \\ 4 & 0 + 16 \cdot 5 + 8 \cdot 4 + 4 \cdot 4 + 2 \cdot 9 & 146 & 22 \\ \end{array}\]\[\begin{array}{c|l|r|r}k & B_k = \sum_j \omega^{jk} b_j & & B_k \bmod 31 \\ \hline 0 & 9 + 8 + 4 + 9 + 3 & 33 & 2 \\ 1 & 9 + 2 \cdot 8 + 4 \cdot 4 + 8 \cdot 9 + 16 \cdot 3 & 161 & 6 \\ 2 & 9 + 4 \cdot 8 + 16 \cdot 4 + 2 \cdot 9 + 8 \cdot 3 & 147 & 23 \\ 3 & 9 + 8 \cdot 8 + 2 \cdot 4 + 16 \cdot 9 + 4 \cdot 3 & 237 & 20 \\ 4 & 9 + 16 \cdot 8 + 8 \cdot 4 + 4 \cdot 9 + 2 \cdot 3 & 211 & 25 \\ \end{array}\]Поточечное произведение $C_k = A_k B_k \bmod 31$: $22 \cdot 2 = 44 \equiv 13$; $16 \cdot 6 = 96 \equiv 3$; $9 \cdot 23 = 207 \equiv 21$; $24 \cdot 20 = 480 \equiv 15$; $22 \cdot 25 = 550 \equiv 23$.\par Обратное преобразование: $c_j = n^{-1} \sum_k C_k \omega^{-jk}$, $\omega^{-1} = \omega^{4} \equiv 16$, $n^{-1} = 5^{-1} \equiv 25 \pmod{31}$ ($5 \cdot 25 = 125 = 4 \cdot 31 + 1$), матрица $(\omega^{-jk} \bmod m)$:\[W^{-} = \begin{pmatrix}1 & 1 & 1 & 1 & 1\\1 & 16 & 8 & 4 & 2\\1 & 8 & 2 & 16 & 4\\1 & 4 & 16 & 2 & 8\\1 & 2 & 4 & 8 & 16\end{pmatrix}.\]\[\begin{array}{c|l|r|r|r}j & \sum_k \omega^{-jk} C_k & & \bmod 31 & c_j = n^{-1} \cdot (\ldots) \\ \hline 0 & 13 + 3 + 21 + 15 + 23 & 75 & 13 & 15 \\ 1 & 13 + 16 \cdot 3 + 8 \cdot 21 + 4 \cdot 15 + 2 \cdot 23 & 335 & 25 & 5 \\ 2 & 13 + 8 \cdot 3 + 2 \cdot 21 + 16 \cdot 15 + 4 \cdot 23 & 411 & 8 & 14 \\ 3 & 13 + 4 \cdot 3 + 16 \cdot 21 + 2 \cdot 15 + 8 \cdot 23 & 575 & 17 & 22 \\ 4 & 13 + 2 \cdot 3 + 4 \cdot 21 + 8 \cdot 15 + 16 \cdot 23 & 591 & 2 & 19 \\ \end{array}\]Проверка: прямое вычисление циклической свёртки над $\mathbb{Z}$, $c_j = \sum_i a_i b_{(j - i) \bmod n}$, даёт $(139, 129, 169, 115, 174)$, по модулю $31$ это $(15, 5, 14, 22, 19)$ --- совпадает. Здесь $\max c_j = 174 \ge 31$: ТЧП даёт свёртку лишь по модулю $m$; для точной свёртки цифр $0, \ldots, 9$ заведомо достаточно модуля $m > n \cdot 9^2 = 405$.\[c = a * b = (15, 5, 14, 22, 19) \pmod{31}.\]
\end{example}
